\documentclass[10pt,a4paper]{amsart}
\usepackage[margin=19mm]{geometry}
\usepackage{amsmath,amssymb,mathtools}
\usepackage[scr=rsfs,cal=euler]{mathalfa}
\usepackage{xfrac}
\usepackage[unicode,hidelinks,bookmarks=false]{hyperref}
\newcommand{\ie}{i.e.\ }
\newcommand{\cf}{cf.\ }
\newcommand{\resp}{respectively\ }
\newcommand{\Subsection}{\subsection}
\providecommand{\nobreakdash}{\nobreak\hbox{-}}
\setlength{\emergencystretch}{2em}
\multlinegap0pt
\usepackage{bm}
\usepackage{bbm}
\usepackage{dsfont}
\usepackage{xspace}
\usepackage{cancel}
\usepackage{mathtools}
\usepackage{amsthm}
\makeatletter
\@ifundefined{theorem}{\newtheorem{theorem}{Theorem}}{}
\@ifundefined{prop}{\newtheorem{prop}[theorem]{Proposition}}{}
\makeatother
\title[Quantum Modular Forms and Resurgence]{Quantum Modular Forms and Resurgence}
\author{Eleanor McSpirit, Larry Rolen}
\date{Source: arXiv:2505.00799v2 (CC BY 4.0)}
\begin{document}
\maketitle

\noindent\emph{Source and licence.} Quantum Modular Forms and Resurgence. Version arXiv:2505.00799v2 (CC BY 4.0), distributed under the Creative Commons Attribution 4.0 International licence, as recorded in the arXiv licence metadata for this exact version and shown on the abstract page https://arxiv.org/abs/2505.00799v2. Full rights notice: licenses/arxiv-2505.00799-CC-BY-4.0.txt.\par\smallskip

\noindent\textit{Editorial scope.} Counted unit: the Prop whose source label is \texttt{prop: modular slash} in the article (source0.tex lines 817--820, source label \texttt{prop: modular slash}), delivered together with the article's own complete original proof of it (source0.tex lines 822--871, one \texttt{proof} environment of 50 lines). No mathematics is written, repaired or re-derived: statement and proof are the article's own text line for line. Required context retained uncounted: prop: med resum (lines 628--629); prop: eichler discont (lines 788--791). Every definition, display and label of those lines is retained unchanged. Omitted from this package and not redistributed: 1--627, 630--787, 792--816, 821--821, 872--1005. Nothing inside a retained excerpt is omitted or altered: every retained body is delivered exactly as the article's lines are written, so no omission receipt of any kind accompanies this package. Citations: the retained excerpts carry no citation command at all, so no bibliography entry of the article is transcribed and no reference is redistributed. Reference adaptation: none was needed and none was made; every reference inside the delivered unit resolves inside it, so no unresolved reference or citation is produced. Printed slips kept literal: no printed word, symbol, hypothesis or number of the counted statement or of its proof was altered. Layout and class: the article's own class is \texttt{amsart} with packages including amssymb, amsthm, cite, hyperref, amsmath, mathtools, amscd, enumitem, caption, verbatim, color, geometry, tikz, setspace; the wrapper is the lane's pinned \texttt{amsart} editorial wrapper at 10pt on A4, which restates the theorem-like environments the retained text needs and adds the article's own macro definitions that the retained text needs (none), with extra wrapper packages (bm, bbm, dsfont, xspace, cancel, mathtools). The delivered numbering is the unit's own. Assets: no retained span contains a figure environment, a graphics-inclusion command, a PDF-page inclusion, a table or a TikZ picture; no image file is copied into this package, the package declares no asset dependency and no image file is redistributed with it. Licence: version arXiv:2505.00799v2 of this article is distributed under the Creative Commons Attribution 4.0 International licence, as recorded in the arXiv licence metadata for this exact version and shown on the abstract page https://arxiv.org/abs/2505.00799v2. A full rights notice is delivered at licenses/arxiv-2505.00799-CC-BY-4.0.txt. Characters: every retained excerpt is pure ASCII, so the delivered engine, XeLaTeX with its default Latin Modern text font, reports no missing character.
\begin{prop}\label{prop: med resum} Suppose $g$ is of Type 4 with respect to $\Gamma_0(N)$ or of Type 5, and suppose $\gamma = \begin{psmallmatrix} a & b \\ c & d\end{psmallmatrix} \in \Gamma_0(N)$. Then $g= \mathcal{S}_{med}^{\pi/2}(\varphi_{-d/c})$.	
\end{prop}

\begin{prop}\label{prop: eichler discont}
	If $g$ is of Type 4 with $2-k= \kappa$ or of Type 5 with $s = \kappa/2$, then  for $\gamma = \begin{psmallmatrix} a & b \\ c & d \end{psmallmatrix} \in \Gamma_0(N)$, we have that 
 $2g|_{\kappa}\gamma(z) = \emph{disc}_\kappa^{\pi/2}(\mathcal{B}_\kappa(\varphi_{-d/c}(z))$.
\end{prop}

\begin{prop} \label{prop: modular slash}
If $g$ is of Type 1 or 2 with $\kappa = k$ or Type 3 with $\kappa = 2-k$, then for $\gamma = \begin{psmallmatrix} a & b \\ c & d \end{psmallmatrix} \in \Gamma_0(N)$, we have that 
\[g|_{\kappa}\gamma(z) = \frac{1}{2\sqrt{z+d/c}} \int_{\Gamma}e^{-u/(z+d/c)}\hat{g}_{-d/c}^{odd}(u)\frac{du}{\sqrt{u}}.\]
\end{prop}

\begin{proof}	
In the cases where $g$ is of \emph{Type 1} or \emph{2}, since the asymptotic series $\varphi_{-d/c}(z)$ is zero, we have that $\mathcal{B}_{1/2}(\varphi_{-d/c}(z))=0$. So, it suffices to analyze 
 \[\frac{1}{\sqrt{u}}\widehat{g}_{-d/c}^{odd}(u) = -\frac{1}{\sqrt{\pi }} \sum_{n=0}^\infty \frac{ L(g, \zeta_c^{-d};-n-1/2)}{\Gamma(2n+2)}(8\pi i)^{n+1/2}u^{n}.
\]
 
We first consider $g$ of \emph{Type 1} or \emph{2}. We assume the notation and some of the techniques of Proposition \ref{prop: med resum}. 
We can write the above as 

 \begin{align*}\frac{1}{\sqrt{u}}\widehat{g}_{-d/c}^{odd}(u) &= - \sum_{n=0}^\infty \frac{ L(g, \zeta_c^{-d};-n-1/2)}{\Gamma(n+1)\Gamma(n+3/2)}(2\pi i)^{n+1/2}u^{n}\\
 &=C_{k, -d/c} \frac{i^{-k}}{\pi} \sum_{n=0}^\infty \Big\lbrack \left(\frac{c}{2\pi}\right)^{k+2n+1}\frac{\Gamma(n+k+1/2)}{\Gamma(n+1)}\\ & \quad \quad \quad \quad\quad \quad\quad \quad\cdot L(g, \zeta_c^{a}; n+k+1/2) (-1)^n(2\pi i)^{n+1/2}u^n\Big\rbrack,
\end{align*}
where $C_{k, -d/c}=1$ when $k$ is an integer and equals $\varepsilon_d^{2k}\left(\frac{c}{d}\right)$ when $k$ is a half-integer. 


If we define
\[\Phi(z) =C_{k, -d/c}\frac{i^{-k}}{\pi} \sum_{n=0}^\infty \left(\frac{c}{2\pi}\right)^{k+2n+1} (-1)^n\Gamma(n+k+1/2)L(g, \zeta_c^{a}; n+k+1/2)(2\pi i)^{n+1/2}z^n,
\]
then we see that $\mathcal{B}_1(\Phi(z)) = \frac{1}{\sqrt{u}}\widehat{g}_{-d/c}^{odd}(u)$. Instead consider $\mathcal{B}_{k+1/2}(\Phi(z))$, which equals 
\[i^{-k} C_{k, -d/c}\sum_{n=0}^\infty \left(\frac{c}{2\pi}\right)^{k+2n+1} (-1)^nL(g, \zeta_c^{a}; n+k+1/2)(2\pi i)^{n+1/2}u^n.
\]
Writing $L(g, \zeta_c^{a}; n+k+1/2)$ as its Dirichlet series and using absolute convergence allows us to write this as 
\begin{align*}&\frac{-C_{k, -d/c}}{\pi i c^k} \sum_{m=1}^\infty \sum_{n=0}^\infty \frac{a(m)\zeta_c^{am}}{(2\pi i m/c^2)^{n+k+1/2}}u^n = \frac{C_{k, -d/c}}{\pi i c^k}\sum_{m=1}^\infty  \frac{a(m)\zeta_c^{am}}{u - 2\pi i m/c^2}\frac{1}{(2\pi i m/c^2)^{k-1/2}}.
\end{align*}
We can argue along the same lines as in Proposition \ref{prop: med resum} in order to move between the Borel-Laplace sum of $\Phi$ with shift 1 and with shift $k+1/2$. We have that
\begin{align*}
&\frac{1}{2\sqrt{z+d/c}} \int_{\Gamma}e^{-u/(z+d/c)}\widehat{g}_{-d/c}^{odd}(u)\frac{du}{\sqrt{u}}\\
&= \frac{1}{2\sqrt{z+d/c}} \int_{\Gamma}e^{-u/(z+d/c)}\mathcal{B}_1(\Phi)\;du\\
%&= \frac{1}{2}\sqrt{z+d/c}(z+d/c)^{-1}\int_\Gamma e^{-u/(z+d/c)}\mathcal{B}_1(\Phi)\;du\\
%&= \frac{1}{2}\sqrt{z+d/c}\cdot (z+d/c)^{-1}\left(\int_0^{e^{i(\frac{\pi}{2}-\varepsilon)}\infty}- \int_0^{e^{i(\frac{\pi}{2}+\varepsilon)}\infty}\right) e^{-u/(z+d/c)}\mathcal{B}_1(\Phi)\;du\\
&= \frac{1}{2}\sqrt{z+d/c}\left( \mathcal{L}_{1}^{\pi/2-\varepsilon}(\mathcal{B}_1(\Phi)) - \mathcal{L}_{1}^{\pi/2+\varepsilon}(\mathcal{B}_1(\Phi))\right) \\
&= \frac{1}{2}\sqrt{z+d/c}\left( \mathcal{L}_{k+1/2}^{\pi/2-\varepsilon}(\mathcal{B}_{k+1/2}(\Phi)) - \mathcal{L}_{k+1/2}^{\pi/2+\varepsilon}(\mathcal{B}_{k+1/2}(\Phi))\right)\\
&= \sqrt{z+d/c}\cdot c^{-k} (z+d/c)^{-k-1/2} \sum_{m=1}^\infty a(m)\zeta_c^{am}\frac{(2\pi i m/c^2)^{k-1/2}}{(2\pi i m/c^2)^{k-1/2}} e^{2\pi i m/c^2(z+d/c)}\\
&= (cz+d)^{-k}\sum_{m=1}^\infty a(m) e^{2\pi i m\left(\frac{az+b}{cz+d}\right)}.
\end{align*}

In the case where $g$ is of \emph{Type 3}, the argument is largely the same. We assume that $g$ is the Eichler integral of $f$. Turning our attention back to 
\[\frac{1}{\sqrt{u}}\widehat{g}_{-d/c}^{odd}(u) = -\frac{1}{\sqrt{\pi }} \sum_{n=0}^\infty \frac{ L(g, \zeta_c^{-d};-n-1/2)}{\Gamma(2n+2)}(8\pi i)^{n+1/2}u^{n},
\]
we see that this equals 
\begin{align*}
	&-\frac{1}{\sqrt{\pi}}\sum_{n=0}^\infty \frac{L(f, \zeta_c^{-d};-n+k-3/2)}{\Gamma(2n+2)}(8\pi i)^{n+1/2} u^{n} \\
	%&= \frac{-i^{-k} \varepsilon_d^{2k} \left( \frac{c}{d}\right)}{\pi}\sum_{n=0}^\infty \frac{(8\pi i n)^{n+1/2}c^{2n-k+3} \sin(\pi(n-k+5/2))}{(2\pi)^{2n-k+3} 2^{2n+1}} \frac{\Gamma(n-k+5/2)}{\Gamma(n+1)} L(f, \zeta_c^{a}; n+3/2) u^n\\
	&= \frac{- c^{k-2} \varepsilon_d^{2k} \left( \frac{c}{d}\right)}{\pi i}\sum_{n=0}^\infty  \frac{\Gamma(n-k+5/2)}{\Gamma(n+1)} \frac{L(f, \zeta_c^{a}; n+3/2)}{(2\pi i/c^2)^{n+5/2-k}} u^n.\\
\end{align*}
Again we define 
\[\Phi(z) := \frac{- c^{k-2} \varepsilon_d^{2k} \left( \frac{c}{d}\right)}{\pi i}\sum_{n=0}^\infty \Gamma(n-k+5/2) \frac{L(f, \zeta_c^{a}; n+3/2)}{(2\pi i/c^2)^{n+5/2-k}} z^n,
\]
so that $\mathcal{B}_1(\Phi(z)) =\frac{1}{\sqrt{u}}\widehat{g}_{-d/c}^{odd}(u)$. After splitting off finitely many terms to ensure convergence of the Dirichlet series representation of $L(f, \zeta_c^{a}; n+3/2)$, we take the Borel transform of the remaining terms of $\Phi$ with shift $5/2-k$. 

We proceed as before to realize the discontinuity after shifting as in Proposition \ref{prop: eichler discont}, deforming contours and relating Borel-Laplace sums as for \emph{Types 1} and \emph{2}. \end{proof}

\end{document}
